\documentclass[tikz,border=2mm]{standalone}
\usepackage{chemfig}
\setchemfig{atom sep=2em}
\begin{document}
\schemestart 	\subscheme{\chemfig{NH_3} 		\arrow{->}[0,0.7] 		\chemfig{R-NH_2} 		\arrow{->}[0,0.7] \chemfig{R-NH-R'} 		\arrow{->}[0,0.7] 		\chemfig{R-N(-[6]R'')-R''}} \schemestop --> ![...](imagenes/tema08/aminas.svg){style="display: block; margin: 0 auto; width: 60%;"} * Se pueden nombrar citando los radicales en orden alfabético seguidos de la terminación **"-amina"**. * En aminas complejas, se puede tomar la cadena más larga como hidrocarburo principal terminando en **"-amina"** y los demás radicales unidos al nitrógeno se nombran antecedidos por la letra **N-**. * Si no actúa como función principal, el grupo amino se nombra como prefijo **"amino-"**. Ejemplos: * **Etilamina**: $\ce{\hspace{0.5cm} CH3-CH2-NH2}$ * **N-metiletilamina**: $\ce{\hspace{0.5cm} CH3-CH2-NH-CH3}$ ### **Amidas** Son compuestos derivados de los ácidos carboxílicos en los que se sustituye el grupo $\ce{-OH}$ por un grupo amino ($\ce{-NH2}$), dando el grupo funcional **amida** ($\ce{-CONH2}$). * Se nombran sustituyendo la terminación **"-oico"** o **"-ico"** del ácido de procedencia por el sufijo **"-amida"**. * Si los átomos de hidrógeno del grupo amino están sustituidos por radicales alquílicos, estas posiciones se indican con el prefijo **N-**. Ejemplos: * **Propanamida**: $\ce{\hspace{0.5cm} CH3-CH2-CONH2}$ * **N-metiletanamida**: $\ce{\hspace{0.5cm} CH3-CONH-CH3}$ ### **Nitrocompuestos** Se pueden considerar derivados de los hidrocarburos en los que se substituye uno o más hidrógenos por el grupo "nitro", $\ce{-NO2}$. Se nombran como sustituyentes del hidrocarburo del que proceden indicando con el prefijo "nitro-" y un número localizador su posición en la cadena carbonada. Las insaturaciones tienen preferencia sobre el grupo nitro: $\ce{CH2 = CH - CH2 - NO2 \hspace{1.5cm}}$ 3-nitro-1-propeno ### **Nitrilos o cianuros** Se caracterizan por tener el grupo funcional **"ciano"** ($\ce{-C \equiv N}$) en un extremo de la cadena. * Se pueden nombrar añadiendo el sufijo **"-nitrilo"** al nombre del hidrocarburo de igual número de carbonos. * También se pueden nombrar como **"cianuro de"** seguido del nombre del radical alquílico con la terminación **"-ilo"**. * Cuando actúa como función secundaria se designa mediante el prefijo **"ciano-"**. Ejemplos: * **Propanonitrilo** (o **cianuro de etilo**): $\ce{\hspace{0.5cm} CH3-CH2-CN}$ ### **Orden de prioridad de grupos funcionales** Cuando en un mismo compuesto existen dos o más grupos funcionales, el orden de prioridad dictado por la IUPAC determina cuál de ellos dará nombre al compuesto como función principal. Los demás grupos se considerarán meros sustituyentes y se nombrarán como prefijos. El orden decreciente de prioridad es el siguiente: | Grupo Funcional | Sufijo (Función Principal) | Prefijo (Función Secundaria) | | :--- | :--- | :--- | | 1. **Ácidos carboxílicos** ($\ce{-COOH}$) | Ácido ...oico | Carboxi- | | 2. **Ésteres** ($\ce{-COO-R}$) | ...ato de ...ilo | Alcoxicarbonil- | | 3. **Amidas** ($\ce{-CONH2}$) | ...amida | Carbamoil- | | 4. **Nitrilos** ($\ce{-CN}$) | ...nitrilo | Ciano- | | 5. **Aldehídos** ($\ce{-CHO}$) | ...al | Formil- (u Oxo-) | | 6. **Cetonas** ($\ce{-CO-}$) | ...ona | Oxo- | | 7. **Alcoholes** ($\ce{-OH}$) | ...ol | Hidroxi- | | 8. **Aminas** ($\ce{-NH2}$) | ...amina | Amino- | | 9. **Éteres** ($\ce{-O-}$) | ...oxi... | Alcoxi- | | 10. **Dobles enlaces** ($\ce{C=C}$) | ...eno | | 11. **Triples enlaces** ($\ce{C \equiv C}$) | ...ino | | 12. **Halógenos** ($\ce{-X}$) e **Inertes/Radicales** ($\ce{-R}$) | Siempre secundarios | ## **4. Isomería** Se llaman **isómeros** a aquellos compuestos diferentes que tienen la misma fórmula molecular pero distintas propiedades físicas o químicas debido a una disposición estructural o espacial diferente. **1. Isomería Estructural o Constitucional** Los átomos están unidos de forma diferente en cada molécula. * **De cadena:** Cambia la distribución de la cadena carbonada (ej. *butano* y *2-metilpropano*). * **De posición:** El mismo grupo funcional cambia de ubicación (ej. *1-propanol* y *2-propanol*). * **De función:** Tienen diferente grupo funcional (ej. *etanol* y *dimetil éter*, ambos $\ce{C2H6O}$). **1. Estereoisomería (Isomería Espacial)** Los compuestos presentan idéntica conectividad molecular pero difieren en la orientación tridimensional de sus átomos. * **Isomería Geométrica (cis/trans):** Propia de compuestos con doble enlace rígido $\ce{C=C}$. El isómero *cis* posee los sustituyentes principales al mismo lado, mientras que el *trans* los posee en lados opuestos. * **Isomería Óptica:** Ocurre ante la presencia de un carbono asimétrico o quiral (unido a 4 sustituyentes distintos). Da lugar a enantiómeros, imágenes especulares no superponibles que desvían la luz polarizada hacia la derecha (dextrógiro) o hacia la izquierda (levógiro). ## **5. Reacciones orgánicas principales** **1. Reacciones de Sustitución** Un átomo o grupo de átomos de la molécula sustrato es sustituido por otro. * **Ejemplo (Halogenación de alcanos):** $$\ce{CH4 + Cl2 ->[luz] CH3Cl + HCl}$$ **2. Reacciones de Adición** Ocurren en moléculas insaturadas (con dobles o triples enlaces). Los átomos del reactivo se añaden rompiendo el enlace pi. * Siguen la **Regla de Markovnikov**: El hidrógeno del reactivo se adiciona mayoritariamente al carbono insaturado que tenga ya más átomos de hidrógeno. * **Ejemplo (Hidratación de alquenos):** $$\ce{CH3-CH=CH2 + H2O ->[H+] CH3-CH(OH)-CH3}$$ *(Producto mayoritario: 2-propanol)* **3. Reacciones de Eliminación** Dos átomos o grupos de átomos se escinden de posiciones adyacentes para dar lugar a una insaturación (doble o triple enlace). * Siguen la **Regla de Saytzeff**: Se elimina preferentemente el hidrógeno del carbono vecino que posea menos hidrógenos, obteniéndose el alqueno más sustituido y estable. * **Ejemplo (Deshidratación de alcoholes):** $$\ce{CH3-CH2-CH(OH)-CH3 ->[H2SO4][\Delta] CH3-CH=CH-CH3 + H2O}$$ *(Producto mayoritario: but-2-eno)* **4. Reacciones de Oxidación y Reducción** * **Oxidación:** Un alcohol primario se oxida a aldehído y este consecutivamente a ácido carboxílico. Un alcohol secundario se oxida a cetona. * **Reducción:** Es el proceso inverso al anterior. **Problemas complementarios (Estilo PAU)** Ejercicio 1 Para cada uno de los siguientes procesos, formula la reacción, indica el nombre de los productos y el tipo de reacción orgánica: * **a) Hidrogenación catalítica del 3-metilbut-1-eno.** $$\ce{CH3-CH(CH3)-CH=CH2 + H2 ->[Pt/Pd] CH3-CH(CH3)-CH2-CH3}$$ *Producto:* 2-metilbutano. *Tipo:* Adición (Reducción). * **b) Deshidratación de 1-butanol con ácido sulfúrico en caliente.** $$\ce{CH3-CH2-CH2-CH2OH ->[H2SO4][\Delta] CH3-CH2-CH=CH2 + H2O}$$ *Producto:* but-1-eno. *Tipo:* Eliminación. * **c) Deshidrohalogenación de 2-bromobutano con hidróxido de potasio en etanol.** $$\ce{CH3-CHBr-CH2-CH3 + KOH ->[etanol] CH3-CH=CH-CH3 + KBr + H2O}$$ *Producto:* but-2-eno (mayoritario por la regla de Saytzeff). *Tipo:* Eliminación. * **d) Obtención de metanol a partir de clorometano.** $$\ce{CH3Cl + NaOH -> CH3OH + NaCl}$$ *Producto:* Metanol. *Tipo:* Sustitución nucleófila.
\end{document}
